\documentclass[11pt]{article}
\usepackage[margin=0.9in]{geometry}
\usepackage{booktabs,longtable,amsmath,amssymb,hyperref,enumitem}
\usepackage[T1]{fontenc}
\usepackage{fancyvrb}
\hypersetup{colorlinks=true,linkcolor=blue,urlcolor=blue}
\title{\textbf{Publishing this: repository, DOI, arXiv}\\
\large Third edition}
\author{Karthikeya \\ \small Accretion Labs}
\date{11 October 2026}
\begin{document}
\maketitle

\section*{What changed, and why there is a third edition}

The first edition had six factual errors. The second edition corrected them and introduced
six more operational ones, which a further review caught. Both sets are listed here so you
do not act on a remembered version of either.

\subsection*{Fixed in the second edition}

\begin{enumerate}[nosep]
\item \textbf{The DOI advice was backwards.} Cite the \emph{version-specific} DOI, not the
      concept DOI. Zenodo says so directly.
\item \textbf{The endorsement mechanism was wrong.} Automatic endorsement needs a claimed
      co-authored paper \emph{and} an institutional email, not ``affiliation or prior
      submission''.
\item \textbf{The endorser's responsibility was understated.} It is more than archive fit.
\item \textbf{``Every journal accepts it'' was unsupported}, and CC~BY does not ``match''
      the MIT code licence --- they govern different artifacts.
\item \textbf{``If it builds locally, arXiv will build it too'' is not true.} Check arXiv's
      own log and preview.
\item \textbf{The moderation predictions were speculation} and are removed.
\end{enumerate}

\subsection*{Fixed in this third edition}

\begin{enumerate}[nosep]
\item \textbf{The secret scan was too narrow.} \texttt{--only-verified} can miss a key you
      have just revoked, which is exactly the key you are looking for.
\item \textbf{I asserted your arXiv account history} rather than telling you to check it.
\item \textbf{Archiving before review was presented as necessary.} It is optional, and if
      you do it, the archive is a review snapshot and not the replication package.
\item \textbf{``The licence should read MIT'' and ``cross-list \texttt{physics.comp-ph} and
      nothing else''} were written as requirements. They are choices.
\item \textbf{Repository facts were asserted, not tested.} There is now a fresh-clone
      checklist.
\item \textbf{A wrong section reference} in the endorser table.
\end{enumerate}

I verified the claims in \S3--\S6 against the official pages, quoted where it matters.
Anything I could not verify is flagged as needing your check --- and \S7 is the set of
things only a fresh clone can settle.

\section*{The recommendation, unchanged}

\textbf{Get the independent code review before you post to arXiv.}

A posting is permanent. v2 replaces v1 for readers, but v1 stays visible and the abstract
page shows the history. Six defects have now been found in this work. Every one was found by
reading code or prose against a claim --- never by running the pipeline, which was green
through all of them. The sixth was in the first edition of this very guide.

\S1--\S4 can all proceed while the review is outstanding. Only the submission itself waits.

\section{Pre-flight}

\subsection*{1.1 Rotate the Materials Project key}

Do this first. The key was on screen in a screenshot. Generate a new one at
\url{https://next-gen.materialsproject.org/api}, replace the contents of
\texttt{C:\textbackslash research\textbackslash mp\_api\_key.txt}, and confirm the old one
stops working.

\subsection*{1.2 Scan for the old key without typing it anywhere}

The first edition told you to paste the old key into a \texttt{git grep}. Don't: that puts
the credential into your shell history, and on Windows into \texttt{PSReadLine}'s saved
history file as well. Use a scanner that looks for credential \emph{shapes} instead:

\begin{Verbatim}[fontsize=\small,frame=single]
# both scan history, not just the working tree
gitleaks detect --source . --redact
trufflehog git file://.            # NOT --only-verified; see below
\end{Verbatim}

The second edition told you to pass \texttt{--only-verified}. That is wrong for this job:
verification works by trying the credential against the provider, so a key you have just
rotated fails verification and can be filtered out of the results --- and a revoked key still
sitting in public history is precisely what you are scanning for. Review the verified,
unverified and unknown findings. Check your installed version's flags, since the syntax has
changed between releases.

GitHub's own secret scanning also runs automatically on public repositories, which is a
reason to check before flipping visibility rather than after. And look in the places a
\texttt{.gitignore} never covered: cached API outputs, notebook checkpoints, screenshots,
and any log you pasted into a document.

A \texttt{.gitignore} entry proves nothing about what was committed in the past. If a
scanner finds the key in history, rewriting history is straightforward on a repository that
has never been public and awkward afterwards --- which is the other reason to do this now.

\subsection*{1.3 What will actually be published}

\begin{itemize}[nosep]
\item \texttt{data/} is ignored, so the two 175\,MB GNoME CSVs stay out.
      \texttt{REPRODUCE.md} carries their URLs and SHA-256 sums.
\item The four cached Materials Project files are committed deliberately --- see \S5 for the
      licence position, which I have now verified.
\item \texttt{review/} holds the referee report and four response documents. Publishing
      them is a strength, but read them once as an outsider first.
\item About 10\,MB without \texttt{data/}.
\end{itemize}

\section{Making the repository public}

\begin{Verbatim}[fontsize=\small,frame=single]
gh repo edit --visibility public --accept-visibility-change-consequences
gh repo edit --description "Chemistry-resolved audit of PBE vs r2SCAN stability verdicts in GNoME"
gh repo edit --add-topic dft --add-topic materials-science --add-topic reproducibility \
             --add-topic r2scan --add-topic gnome --add-topic convex-hull
gh api repos/:owner/:repo --jq '.visibility, .has_issues, .license.spdx_id'
gh run list --limit 3
\end{Verbatim}

Issues should be on, since \texttt{REVIEWER.md} asks for findings there. The licence field
should read whatever code licence you have actually chosen --- it is \texttt{MIT} in the
repository today, and confirming that is your decision rather than a publication
requirement. Watch the next CI run: going public changes the Actions quota and
how secrets behave in forks, so confirm the workflow still \emph{skips} the credentialed
phases cleanly rather than failing when \texttt{MP\_API\_KEY} is unavailable.

\section{The DOI --- cite the version, not the concept}

This is the correction that matters most, because getting it wrong makes the replication
package ambiguous. Zenodo's own guidance:

\begin{quote}
``You should normally always use the DOI for the \textbf{specific version} of your record in
citations.''
\end{quote}

A version DOI identifies one exact release. A concept DOI covers all versions and resolves
to the latest, which is useful for pointing at an evolving project and wrong for a
reproducibility citation --- a reader following it in a year gets code that is not the code
the paper used.

\begin{enumerate}[nosep]
\item Sign in to \url{https://zenodo.org} with GitHub. Zenodo only sees public
      repositories, so \S2 comes first.
\item Enable \texttt{FOUNDERACCRETION/gnome-audit} in Zenodo's GitHub settings. It archives
      releases created \emph{after} the hook is on, so enable before releasing.
\item Cut the release. \texttt{.zenodo.json} is in the repository, so the metadata comes
      from there rather than a web form:
\begin{Verbatim}[fontsize=\small,frame=single]
gh release create v1.0.0 \
  --title "v1.0.0 — accompanies the preprint" \
  --notes "Analysis code, pre-registered decision rules and derived results.
53 headline checks and 242 table-value checks green. See REPRODUCE.md and REVIEWER.md."
\end{Verbatim}
\item Zenodo mints both DOIs. \textbf{Put the version DOI and the release tag in the
      paper}; the concept DOI is fine for a README badge.
\item Edit three places: the paper's Data and code availability section (it currently says
      the deposit has not been made), \texttt{CITATION.cff}'s \texttt{doi:} field, and the
      README badge. Then rebuild the PDF and re-run \texttt{python run\_all.py}.
\end{enumerate}

\textbf{Archiving before the review is optional.} A reviewer needs a fixed thing to look
at, and a tagged commit does that job --- \texttt{git tag review-2026-10-11 \&\& git push
--tags} is enough, and costs nothing if the review changes the code. Two workable routes:

\begin{itemize}[nosep]
\item \emph{Tag only.} Send the reviewer a tag, archive once after their findings are
      addressed, and cite that version DOI. Simplest, and the DOI matches the submitted
      paper by construction.
\item \emph{Archive now as a snapshot.} If you want a citable object immediately, archive it
      and \textbf{label it a review snapshot in the release notes} --- not the replication
      package. Then archive again after the review and cite the later version DOI.
\end{itemize}

Either way the rule is the same: the version DOI in the paper must correspond to the code
and outputs that produced the paper's numbers. Do not cite a snapshot as the replication
package for a paper whose numbers came from later code.

\section{arXiv}

\subsection*{4.1 Endorsement: the mechanism, corrected}

arXiv's automatic route requires \emph{both} of:

\begin{itemize}[nosep]
\item that you have ``claimed ownership of a paper submitted by a co-author'', and
\item that your email address meets arXiv's institutional email criteria.
\end{itemize}

arXiv says you ``may receive'' automatic endorsement when both hold, and does not publish
the email criteria. The second edition stated flatly that you have no co-authored arXiv
paper; I have not checked your account, so treat it as a conditional: \emph{if} you have
never claimed authorship of a co-authored arXiv submission, the first condition fails and
the second is moot. Either way, \textbf{do not act on a prediction}. Create the account,
start a submission, and follow the instructions arXiv gives \emph{your} account. That is
faster and more reliable than reasoning about the policy.

\subsection*{4.2 What you are actually asking an endorser for}

The first edition called it ``archive-fit endorsement''. That understates it. arXiv's
wording:

\begin{quote}
``We ask that you only endorse individuals when you are confident they will submit
high-quality research papers.''
\end{quote}

Endorsers should personally know the author or have seen the paper. They should decline if
the author ``is unfamiliar with the basic facts of the field'' or if ``the work is entirely
disconnected with current work in the area''. They are not expected to read the paper in
detail or verify its correctness --- but it is a judgement about you, not a formality.

Mechanically: start a submission, arXiv issues a six-character endorsement code, you send
the code and the paper, the endorser enters it with a yes or no. The ``Which of these authors
are endorsers?'' link on any abstract page shows who is eligible. \textbf{Ask one person at a
time}; arXiv explicitly says not to email several at once or the same person twice.

A request that matches what arXiv is actually asking:

\begin{quote}
Would you consider endorsing me for a \texttt{cond-mat.mtrl-sci} submission? I understand
endorsement is not peer review, so I am sending the manuscript and the replication package
so you can judge whether it is suitable work for the archive. [One sentence on what the
paper does and why it touches their work.] My endorsement code is XXXXXX. I am an
independent researcher at Accretion Labs.
\end{quote}

Two notes on that. Do not describe yourself as ``unaffiliated'' while listing Accretion Labs
in the author block --- pick one description and make it accurate. And the people below are
\textbf{candidates I have not verified}: before approaching anyone, check their current
affiliation and confirm eligibility through the endorser link on one of their arXiv abstract
pages.

\begin{center}
\begin{tabular}{llp{0.36\textwidth}}
\toprule
Candidate & Base (verify) & Why they are a sensible ask \\
\midrule
G.\ Pizzi & PSI / EPFL & Reents--Bercx--Pizzi is the closest prior framing, as your related-work section on flip-rate framing says \\
G.\ Cs\'anyi & Cambridge & Warford--Thiemann--Cs\'anyi's selective-$U$ argument is what Section 4.3 tests \\
G.-M.\ Rignanese & UCLouvain & Yu et al.'s per-element error map is your methodological precedent \\
C.\ J.\ Bartel & Minnesota & decomposition-reaction framing is upstream of your outcome definition \\
Z.\ Wang & UC San Diego & An et al.\ ran the SCAN re-screen your Section 2 compares against \\
\bottomrule
\end{tabular}
\end{center}

\subsection*{4.3 What to upload}

Source, not PDF --- and this is a requirement, not a preference: arXiv does not accept a PDF
generated from TeX source in place of that source.

\begin{Verbatim}[fontsize=\small,frame=single]
cd paper
tar czf arxiv_submission.tar.gz gnome_audit_paper.tex \
    fig1_flip_by_chemistry.png fig2_budget.png fig3_gnome_vs_known.png \
    fig4_per_element.png fig5_bandgap.png fig6_logit.png
\end{Verbatim}

\begin{itemize}[nosep]
\item No \texttt{.bib} is needed: the references are an inline \texttt{thebibliography}, so
      nothing has to run BibTeX.
\item \texttt{\textbackslash graphicspath\{\{./\}\{../results/\}\}} resolves via the
      \texttt{./} branch inside the tarball; the second branch is harmless.
\item Extract the tarball into an empty directory and compile \emph{that}, not your working
      copy, so a file you forgot to include shows up now.
\item Then --- and this is the part the first edition got wrong --- \textbf{read arXiv's own
      compilation log and look at its PDF preview before submitting}. A local build is
      evidence, not a guarantee: compiler version, package availability, filename handling
      and configuration all differ. Nor does source submission guarantee a usable HTML
      rendering; treat that as a bonus if it appears.
\end{itemize}

\subsection*{4.4 Metadata}

\begin{description}[nosep,leftmargin=2.4em]
\item[Primary] \texttt{cond-mat.mtrl-sci}.
\item[Cross-list] \texttt{physics.comp-ph} is the natural second home given how much of
      the paper is method and reproducibility, but it is a judgement about your paper's
      scope rather than a rule. Pick deliberately and keep the list short; scattering
      cross-lists does not help.
\item[Abstract] Plain text with the light TeX arXiv accepts (\texttt{\$\^{}2\$},
      \texttt{\$U\$}, \texttt{\$d\$}). Check the preview; a broken abstract is the most
      visible possible typo.
\item[Comments] \emph{49 pages, 6 figures. Code, pre-registered decision rules and
      derived results:}\\
      \texttt{github.com/FOUNDERACCRETION/gnome-audit}\\
      \emph{archived at doi:…\ (v1.x). Includes an independent-review packet.}
\item[Affiliation] ``Accretion Labs''.
\item[ORCID] Get one at \url{https://orcid.org} and link it. Free, and required for the
      ChemRxiv route anyway.
\end{description}

\subsection*{4.5 Licence}

The licence on a posted version is irrevocable. arXiv offers its own non-exclusive
distribution licence and a set of Creative Commons options including CC~BY~4.0 and
CC~BY-SA~4.0. Different \emph{versions} of a paper can carry different licences, but you
cannot change the one already posted.

Two corrections to the first edition. I claimed the non-exclusive licence is accepted by
``every journal'': I have no basis for that and arXiv itself tells you to check journal and
funder policies. And CC~BY does not ``match'' the MIT licence on the code --- the manuscript
and the code are separate artifacts under separate licences, and choosing one says nothing
about the other. The third artifact, the Materials Project extracts, is governed by neither
(see \S5).

So: decide the manuscript licence after checking the policy of whatever journal you might
submit to later. If you have no specific journal in mind, CC~BY~4.0 is the choice consistent
with a paper whose point is that others can check it. If you do, read that journal's policy
first --- this is the one decision here you cannot walk back.

\subsection*{4.6 After submitting}

Submissions go through moderation for topical and scholarly suitability, and endorsement
does not remove that step: arXiv can still reject or reclassify. A reclassification
suggestion is routine and not a verdict on the work.

The first edition also predicted that a sole author with a company affiliation would get
closer scrutiny, and that the preregistration and AI disclosure were ``exactly what a
moderator wants to see''. I have no evidence for either and both are now removed. Submit the
work and let the process run.

\section{The Materials Project extracts --- verified}

The first edition asserted the licence position. I have now checked it. The Materials
Project terms of use place downloaded content under CC~BY~4.0 and state that it

\begin{quote}
``may be copied, distributed, transmitted, and adapted, without obtaining specific
permission, provided proper attribution is given to the Materials Project''
\end{quote}

The same terms permit collecting large fractions of the database through the API and
presenting processed results with attribution, and prohibit scraping the website directly.
The four committed files came through the API, so redistributing them with attribution is
within the terms. \texttt{results/PROVENANCE.md} now quotes this, names the required
attribution, and states explicitly that \textbf{the repository's MIT licence cannot grant
rights to third-party data} --- the MP extracts are governed by the MP terms and the GNoME
CSVs, which are not redistributed here, by whatever DeepMind applies.

Still worth your own check before publishing: that no credential or private value sits
inside those cached CSVs (they should contain only \texttt{material\_id},
\texttt{energy\_above\_hull}, \texttt{nelements} and \texttt{chemsys}), and that the README
distinguishes cached API pulls from regenerated results. It does.

\section{ChemRxiv, the fallback}

Verified: an ORCID iD is required to register and upload, submissions are PDF with
supplementary material in any format, and the licence choices offered are CC~BY~4.0,
CC~BY-NC~4.0 and CC~BY-NC-ND~4.0. Content must be original research, in English, not already
accepted for publication, and it passes a plagiarism check. ChemRxiv's author FAQ puts
screening at one to two business days and posting one to two days after that.

Two corrections. ORCID is the \emph{registration} gate, not the only gate --- ChemRxiv
curates and screens submissions, and approval does not imply endorsement by its publisher or
partner societies. And treat the timings and licence list as current-at-time-of-writing
details to confirm against its submission guide rather than as guarantees.

It is a real venue, not a consolation prize. The honest trade-off is audience: your paper
lands under computational chemistry rather than among the \texttt{cond-mat.mtrl-sci} readers
who would actually use the budget rule. Try arXiv first; use ChemRxiv if the endorsement
route stalls for two or three weeks.

\section{Order of operations}

\begin{enumerate}[nosep]
\item Rotate the Materials Project key. \hfill (\S1.1, today)
\item Scan history with a secret scanner --- never by typing the key. \hfill (\S1.2)
\item Read \texttt{review/} once as an outsider. \hfill (\S1.3)
\item Tag a review commit and \textbf{send it with \texttt{REVIEWER.md} to an independent
      reviewer.} \hfill (the long pole --- start now)
\item Make the repository public; check Actions, issues, licence. \hfill (\S2)
\item Verify the four fresh-clone facts. \hfill (\S7)
\item Optionally enable Zenodo and archive a labelled review snapshot. \hfill (\S3)
\item Get an ORCID; create an arXiv account; start a submission to see what your account is
      actually told about endorsement. \hfill (\S4.1)
\item Ask one endorser. Wait. Ask the next if needed. \hfill (\S4.2)
\item \textbf{Wait for the review. Fix what it finds.} \hfill (the hold point)
\item If anything scientific changed, cut and archive a new version; put \emph{that} version
      DOI in the paper; rebuild; re-run \texttt{run\_all.py}. \hfill (\S3)
\item Decide the manuscript licence against a real journal policy. \hfill (\S4.5)
\item Upload the source tarball; read arXiv's compilation log and preview; submit.
      \hfill (\S4.3)
\item Add the arXiv ID to \texttt{README.md} and \texttt{CITATION.cff}.
\end{enumerate}

Steps 1--8 run in parallel with step 4. Step 9 blocks steps 10--12, and step 10 exists
because of step 9.

\section{Verify these from a fresh clone, not from this guide}

Everything this guide says about the repository is a reported project fact. I have looked at
the working tree on your machine, which is not the same as what a clone gets --- a
\texttt{.gitignore} change, for instance, does nothing about what is already tracked or
untracked until you commit. Before publishing, run this from a clone in an empty directory:

\begin{Verbatim}[fontsize=\small,frame=single]
git clone https://github.com/FOUNDERACCRETION/gnome-audit fresh && cd fresh
# 1. are the four cached Materials Project files actually tracked?
git ls-files results/ | grep -E 'mp_gga_raw|mp_r2scan_raw|known_baseline|mp_completeness'
# 2. is anything that should be ignored present, or anything needed missing?
ls results/ | wc -l ; cat .gitignore
# 3. does a no-credential run work, and does it label stored values as stored?
#    (download the two CSVs into data/ first, per REPRODUCE.md section 2)
find results -type f \( -name '*.csv' -o -name '*.png' -o -name '*.json' \) \
     ! -name 'known_baseline.csv' ! -name 'mp_completeness.csv' \
     ! -name 'mp_gga_raw.csv' ! -name 'mp_r2scan_raw.csv' -delete
unset MP_API_KEY MP_API_KEY_FILE
python run_all.py 2>&1 | tee /tmp/fresh.log
grep -c 'stored, not regenerated here' /tmp/fresh.log   # must be non-zero
grep -i 'api.key\|credential' /tmp/fresh.log            # must find no demand for one
\end{Verbatim}

Four things that must hold, and none of which this guide can establish for you:

\begin{itemize}[nosep]
\item the four cached inputs are tracked, so a clone has them;
\item the wipe in step 3 preserves exactly those four and nothing else;
\item the advertised reproducible path needs no credential, and says which values are
      stored rather than regenerated;
\item \texttt{results/PROVENANCE.md} and the licence files are present in the clone.
\end{itemize}

Then confirm the release metadata matches the final version: the title, the version DOI, and
the tag the paper cites.

\section{What the repository carries for a reviewer}

\begin{itemize}[nosep]
\item \texttt{REVIEWER.md} --- reproduction versus correctness; all defects found so far and
      how each was found, as calibration for how much the green checks are worth; eight
      ordered areas to inspect with the specific thing to check in each; and a request to
      independently reimplement the headline cell rates and the budget comparison rather
      than rerun ours.
\item \texttt{REPRODUCE.md} --- built around the step that makes the rest meaningful
      (delete the stored results first), with input checksums and a troubleshooting order.
\item \texttt{results/runtimes.csv} --- per-phase wall time, total, platform, Python version,
      written on every run.
\item \texttt{results/PROVENANCE.md} --- what is regenerated, what is cached, the verified
      Materials Project terms, and the limits of the MIT licence.
\item \texttt{CITATION.cff}, \texttt{.zenodo.json}.
\item \texttt{NB} and \texttt{SEED} environment variables on the bootstrap phases, so a
      reviewer can rerun a stability check cheaply or expensively without editing code.
\end{itemize}

\end{document}
